%% Drafted from the mapped corpus by scripts/draft_topics.py.
% Model: gpt-5.6-terra; source characters: 19574.
\begin{konspekt}
\kreq{Дефиниции}{контекстносвободна граматика, извод, контекстносвободен език, дърво на извод --- \kref{5.1}.}
\kreq{Конструкции, без доказателства}{затвореност относно регулярните операции --- обединение, конкатенация, звезда на Клини --- \kref{5.4}.}
\kreq{Формулировка и доказателство}{лема за разрастването за КСЕ --- \kref{5.2}.}
\kreq{Пример с доказателство}{език, който не е контекстносвободен --- \kref{5.3}.}
\kreq{Доказателство}{незатвореност относно сечение --- \kref{5.5}.}
\kreq{Доказателство}{незатвореност относно допълнение --- \kref{5.6}.}
\kreq{Типове задачи}{по даден КСЕ да се построи КСГ с обосновка; да се определи с обосновка дали даден език е контекстносвободен.}
\kreq{Литература}{[34], [38], [43].}
\end{konspekt}

\section{Контекстносвободни граматики и дървета на извода}\label{s:5.1}

\begin{defn}
Контекстносвободна граматика (КСГ) е четворка
\[
G=\langle V,\Sigma,R,S\rangle,
\]
където $V$ е крайно множество от нетерминали, $\Sigma$ е крайна азбука от терминали, $V\cap\Sigma=\varnothing$, $S\in V$ е начален нетерминал и
\[
R\subseteq V\times(V\cup\Sigma)^*
\]
е крайно множество от правила. Правилото $(A,\alpha)\in R$ се записва като $A\to\alpha$.
\end{defn}

\begin{defn}[Едностъпков и краен извод]
Едностъпковият извод, породен от $G$, се означава с $\Rightarrow_G$. По определение
\[
\lambda A\rho\Rightarrow_G\lambda\alpha\rho,
\]
когато $A\to\alpha$ е правило на $G$ и $\lambda,\rho\in(V\cup\Sigma)^*$.
С $\Rightarrow_G^*$ означаваме рефлексивно-транзитивното затваряне на тази релация.
\end{defn}

\begin{defn}
Езикът, породен от граматиката $G$, е
\[
L(G)=\{\omega\in\Sigma^*\mid S\Rightarrow_G^*\omega\}.
\]
Език $L\subseteq\Sigma^*$ се нарича \emph{контекстносвободен език} (КСЕ), ако съществува КСГ $G$, за която $L=L(G)$.
\end{defn}

\begin{defn}[Дърво на извод]
Дърво на извод в $G$ е крайно наредено кореново дърво. Коренът е означен с начален нетерминал. Ако вътрешен връх е означен с $A$, последователните му деца са означени със символите от дясната страна на някое правило $A\to\alpha$. При $A\to\varepsilon$ използваме едно листо $\varepsilon$. Резултатът е редицата от означенията на листата отляво надясно, като $\varepsilon$ не добавя символ. За завършено дърво всички листа са терминали или $\varepsilon$.
\end{defn}

\begin{defn}[Най-ляв извод]
Извод е най-ляв, когато във всяка стъпка се заменя най-левият нетерминал в текущата дума.
\end{defn}

\begin{prop}[Изводи и дървета] За дума $\omega\in\Sigma^*$ са еквивалентни следните твърдения:

\begin{enumerate}
\item $S\Rightarrow_G^*\omega$, тоест $\omega\in L(G)$;
\item съществува дърво на извод с корен $S$ и резултат $\omega$;
\item съществува най-ляв извод на $\omega$ от $S$.
\end{enumerate}

\end{prop}
\begin{proof}
От извод строим дървото по индукция по броя стъпки: началното дърво е един лист $S$, а замяната на нетерминал разширява съответния лист с децата от приложеното правило. Резултатът на частичното дърво е текущата дума. Обратно, за дадено завършено дърво разширяваме винаги най-левия останал нетерминален лист според неговите деца в дървото. Крайността гарантира завършване и получаваме най-ляв извод на резултата. Всеки най-ляв извод е частен случай на извод. Това доказва трите импликации.
\end{proof}

\section{Лема за разрастването}\label{s:5.2}

Лемата за разрастването, наричана още лема за покачването, дава необходимо условие един език да е контекстносвободен. Тя е особено полезна за доказване, че даден език \emph{не} е контекстносвободен.

\begin{keythm}[Лема за разрастването за контекстносвободни езици]
Нека $L\subseteq\Sigma^*$ е контекстносвободен език. Тогава съществува число $p>0$, наречено \emph{константа на разрастването}, такова че за всяка дума $\alpha\in L$ с $|\alpha|\geq p$ съществува разлагане
\[
\alpha=xyuvw,
\]
за което са изпълнени:

\begin{enumerate}
\item $|yv|\geq 1$;
\item $|yuv|\leq p$;
\item за всяко $i\in\mathbb{N}$ е изпълнено
\[
xy^iuv^iw\in L.
\]
\end{enumerate}
\end{keythm}

Идеята е, че достатъчно дълга дума има достатъчно високо дърво на извод. По един път в това дърво неизбежно се повтаря нетерминал. Частта от дървото между двете срещания на същия нетерминал може да се повтаря произволен брой пъти или да се премахва.

\begin{proof}
Нека $G=\langle V,\Sigma,R,S\rangle$ е КСГ, за която $L=L(G)$. Без ограничение на общността приемаме, че $G$ е в нормална форма на Чомски. За непразните думи, които разглеждаме, всяко използвано правило е от един от видовете
\[
A\to BC
\qquad\text{или}\qquad
A\to a,
\]
където $A,B,C\in V$ и $a\in\Sigma$.

Нека
\[
k=|V|
\qquad\text{и}\qquad
p=2^k.
\]
Ще покажем, че $p$ е подходяща константа.

Използваме следното елементарно наблюдение за двоични дървета: ако във всяко корен--листо пътче има по-малко от $k$ ребра, то дървото има по-малко от $2^k$ листа. Действително, на всяко ниво броят на върховете може най-много да се удвои, а листата са не повече от броя на възможните пътища с дължина, по-малка от $k$.

Нека сега $\alpha\in L$ и $|\alpha|\geq p=2^k$. Разглеждаме дърво на извод за $\alpha$ и в него избираме най-дълъг корен--листо път. Тъй като граматиката е в нормална форма на Чомски, дървото е двоично: всеки вътрешен връх има най-много две деца. Броят на листата е $|\alpha|$, следователно дървото има поне $2^k$ листа. Ако всеки път имаше най-много $k$ нетерминални върха, щеше да има най-много $k-1$ двоични разклонявания и най-много $2^{k-1}$ терминални листа. Понеже листата са поне $2^k$, избраният най-дълъг път има поне $k+1$ нетерминални върха.

По този път се срещат поне $k+1$ нетерминала, а нетерминалите в $V$ са само $k$. Следователно, по принципа на Дирихле, по пътя има два върха, означени с един и същ нетерминал $A$.

Сред последните $k+1$ срещания на нетерминали по този най-дълъг път избираме две срещания на един и същ нетерминал $A$. Съответстващата част от дървото дава разлагане
\[
\alpha=xyuvw
\]
и изводи
\[
S\Rightarrow_G^* xAw,
\qquad
A\Rightarrow_G^* yAv,
\qquad
A\Rightarrow_G^* u.
\]
Оттук за всяко $i\in\mathbb{N}$ получаваме
\[
A\Rightarrow_G^* y^iAv^i
\Rightarrow_G^* y^iuv^i,
\]
а следователно
\[
S\Rightarrow_G^* xy^iuv^iw.
\]
Значи
\[
xy^iuv^iw\in L
\]
за всяко $i\in\mathbb{N}$.

Остава да проверим първите две условия. Изводът $A\Rightarrow_G^* yAv$ между двете различни срещания на $A$ съдържа поне една стъпка. При първото разклоняване по избрания път другото дете е корен на поддърво, което в нормална форма на Чомски извежда поне един терминален символ. Този символ попада в $y$ или във $v$. Следователно
\[
|yv|\geq 1.
\]

Горното срещане на $A$ е сред последните $k+1$ нетерминални върха на избрания най-дълъг път. Затова от него до листото има най-много $k$ разклонявания. Понеже избраният път е най-дълъг в цялото дърво, никой страничен път в поддървото на това срещане не е по-дълъг. Следователно височината на поддървото е най-много $k$. То е двоично, затова има най-много $2^k=p$ терминални листа. Тези листа дават точно думата $yuv$, откъдето
\[
|yuv|\leq p.
\]
Така трите условия са доказани.
\end{proof}

\begin{remark}
В лемата е съществено разлагането да съществува за всяка достатъчно дълга дума, но то не е предварително зададено. При приложение за доказване на неконтекстносвободност трябва да се разгледат \emph{всички} разлагания, удовлетворяващи условията на лемата, и за всяко от тях да се намери стойност на $i$, при която получената дума не принадлежи на разглеждания език.
\end{remark}

\section{Приложение: език, който не е контекстносвободен}\label{s:5.3}

\begin{example}
Езикът
\[
L=\{a^n b^n c^n\mid n\in\mathbb{N}\}
\]
не е контекстносвободен.
\end{example}

\begin{proof}
Да допуснем противното: нека $L$ е контекстносвободен. Нека $p$ е константата от лемата за разрастването. Избираме думата
\[
\alpha=a^p b^p c^p\in L.
\]
Тя има дължина поне $p$, следователно според лемата има разлагане
\[
\alpha=xyuvw,
\]
за което
\[
|yuv|\leq p,
\qquad
|yv|\geq 1,
\]
и
\[
xy^iuv^iw\in L
\]
за всяко $i\in\mathbb{N}$.

Тъй като всяка от трите еднородни части $a^p$, $b^p$, $c^p$ има дължина $p$, поддумата $yuv$ с дължина най-много $p$ не може да съдържа едновременно буквите $a$, $b$ и $c$. Следователно $y$ и $v$ засягат най-много две съседни групи от буквите.

Разглеждаме думата при $i=2$:
\[
xy^2uv^2w.
\]
Поне едно от $y,v$ е непразно.

Ако $y$ и $v$ са съставени само от една буква, повторението увеличава броя на поне една от буквите $a,b,c$, без да увеличава броя и на трите по един и същ начин. Следователно в получената дума броят на $a$, $b$ и $c$ не е еднакъв, така че тя не принадлежи на $L$.

Ако някое от $y$ и $v$ съдържа две различни букви, то то пресича границата между две от групите $a^p$, $b^p$, $c^p$. При повторяването му се нарушава формата
\[
a^*b^*c^*.
\]
Следователно и в този случай
\[
xy^2uv^2w\notin L.
\]

Получаваме противоречие с третото условие на лемата за разрастването. Следователно $L$ не е контекстносвободен.
\end{proof}

\section{Затвореност относно регулярните операции}\label{s:5.4}

Конспектът иска \emph{конструкции} за затвореността относно регулярните
операции --- обединение, конкатенация и звезда на Клини --- без доказателства.
Конструкциите по-долу работят за произволни входни граматики, а не само за
конкретните примери от следващите раздели.

Нека
\[
G_1=\langle V_1,\Sigma,R_1,S_1\rangle
\qquad\text{и}\qquad
G_2=\langle V_2,\Sigma,R_2,S_2\rangle
\]
са контекстносвободни граматики над обща азбука $\Sigma$. Преди всяка от
конструкциите преименуваме нетерминалите така, че
\[
V_1\cap V_2=\varnothing,
\]
и избираме нов начален нетерминал $S\notin V_1\cup V_2$. Преименуването не
променя породения език: то е просто биекция върху множеството от нетерминали,
приложена едновременно към лявата и дясната страна на всяко правило. Без него
двете граматики биха си споделили нетерминал и изводите им биха се смесили.

\begin{prop}[Конструкция за обединение]
Езикът $L(G_1)\cup L(G_2)$ се поражда от граматиката
\[
G_\cup=\langle V_1\cup V_2\cup\{S\},\ \Sigma,\ R_1\cup R_2\cup\{S\to S_1,\ S\to S_2\},\ S\rangle .
\]
\end{prop}
\begin{proof}
Първата стъпка на всеки извод от $S$ е $S\to S_1$ или $S\to S_2$; други правила
за $S$ няма. Тъй като $V_1\cap V_2=\varnothing$, след избора на $S_1$ всички
достижими нетерминали са от $V_1$ и приложимите правила са точно тези на $R_1$,
тоест изводът продължава като извод в $G_1$; аналогично за $S_2$ и $G_2$.
Следователно $S\Rightarrow^*_{G_\cup}\omega$ точно когато $S_1\Rightarrow^*_{G_1}\omega$
или $S_2\Rightarrow^*_{G_2}\omega$.
\end{proof}

\begin{prop}[Конструкция за конкатенация]
Езикът $L(G_1)L(G_2)=\{\omega_1\omega_2\mid \omega_1\in L(G_1),\ \omega_2\in L(G_2)\}$
се поражда от граматиката
\[
G_\cdot=\langle V_1\cup V_2\cup\{S\},\ \Sigma,\ R_1\cup R_2\cup\{S\to S_1S_2\},\ S\rangle .
\]
\end{prop}
\begin{proof}
Единственото правило за $S$ дава $S\Rightarrow S_1S_2$. Двата нетерминала
разполагат с непресичащи се множества нетерминали, затова поддървото над $S_1$
използва само правила от $R_1$, а поддървото над $S_2$ --- само правила от $R_2$.
Резултатът е конкатенация на дума, изводима от $S_1$ в $G_1$, и дума, изводима
от $S_2$ в $G_2$; обратно, всяка такава двойка дава извод в $G_\cdot$.
\end{proof}

\begin{prop}[Конструкция за звезда на Клини]
Езикът $L(G_1)^*$ се поражда от граматиката
\[
G_*=\langle V_1\cup\{S\},\ \Sigma,\ R_1\cup\{S\to S_1S,\ S\to\varepsilon\},\ S\rangle .
\]
\end{prop}
\begin{proof}
Правилата за $S$ пораждат точно думите $S_1^k$ за $k\in\mathbb{N}$: $k$ пъти се
прилага $S\to S_1S$ и веднъж $S\to\varepsilon$. Всяко от $k$-те срещания на
$S_1$ извежда независимо дума от $L(G_1)$, тъй като $S\notin V_1$ и правилата на
$R_1$ не могат да въведат $S$. При $k=0$ получаваме $\varepsilon$, което
принадлежи на $L(G_1)^*$ по определение.
\end{proof}

\begin{remark}
Трите конструкции обясняват и защо аналогичен подход не работи за сечение: няма
как с едно правило да се наложи \emph{една и съща} дума да бъде изведена
едновременно в двете граматики. Следващите два раздела показват, че такава
конструкция изобщо не съществува.
\end{remark}

\section{Незатвореност относно сечение}\label{s:5.5}

Класът на контекстносвободните езици е затворен относно обединение, конкатенация и звездата на Клини --- това дадоха конструкциите от предходния раздел. В частност, ако $L_1$ и $L_2$ са КСЕ над една и съща азбука, то $L_1\cup L_2$ е КСЕ. За сечение обаче съответното свойство не е вярно.

\begin{thm}
Класът на контекстносвободните езици не е затворен относно сечение.
\end{thm}

\begin{proof}
Разглеждаме езиците
\[
L_1=\{a^n b^n c^k\mid n,k\in\mathbb{N}\}
\]
и
\[
L_2=\{a^n b^k c^k\mid n,k\in\mathbb{N}\}.
\]
И двата езика са контекстносвободни. Например за $L_1$ може да се използва граматиката с правила
\[
S\to AC,\qquad
A\to aAb\mid\varepsilon,\qquad
C\to cC\mid\varepsilon.
\]
Нетерминалът $A$ поражда еднакъв брой букви $a$ и $b$, а $C$ поражда произволен брой букви $c$.

Аналогично за $L_2$ може да се използва граматиката
\[
S\to AD,\qquad
A\to aA\mid\varepsilon,\qquad
D\to bDc\mid\varepsilon.
\]
Тук $A$ поражда произволен брой букви $a$, а $D$ --- еднакъв брой букви $b$ и $c$.

В дума от сечението $L_1\cap L_2$ броят на $a$ и $b$ трябва да е еднакъв поради принадлежността към $L_1$, а броят на $b$ и $c$ трябва да е еднакъв поради принадлежността към $L_2$. Следователно
\[
L_1\cap L_2
=
\{a^n b^n c^n\mid n\in\mathbb{N}\}.
\]
Но този език не е контекстносвободен според предходния пример. Следователно сечението на два контекстносвободни езика не е непременно контекстносвободно.
\end{proof}

\section{Незатвореност относно допълнение}\label{s:5.6}

Допълнението на език $L\subseteq\Sigma^*$ винаги се разбира относно фиксираната азбука $\Sigma$:
\[
\overline{L}=\Sigma^*\setminus L.
\]

\begin{thm}
Класът на контекстносвободните езици не е затворен относно допълнение.
\end{thm}

\begin{proof}
Да допуснем противното: нека за всеки контекстносвободен език $L$ и неговото допълнение $\overline{L}$ също е контекстносвободно.

Нека $L_1$ и $L_2$ са контекстносвободните езици от доказателството за незатвореност относно сечение. По допускане $\overline{L_1}$ и $\overline{L_2}$ са КСЕ. Тъй като КСЕ са затворени относно обединение,
\[
\overline{L_1}\cup\overline{L_2}
\]
е КСЕ. Прилагаме отново допускането за затвореност относно допълнение и получаваме, че
\[
\overline{\overline{L_1}\cup\overline{L_2}}
\]
също е КСЕ.

По закона на де Морган обаче
\[
\overline{\overline{L_1}\cup\overline{L_2}}
=
L_1\cap L_2.
\]
Това противоречи на вече доказаната незатвореност относно сечение. Следователно класът на контекстносвободните езици не е затворен относно допълнение.
\end{proof}

\begin{supp}[Бележка]
От незатвореността относно сечение сама по себе си не следва непосредствено незатвореност относно допълнение. В доказателството е необходимо и известното свойство, че контекстносвободните езици са затворени относно обединение, за да се приложи законът на де Морган.
\end{supp}

\section{Изпитен фокус}\label{s:5.7}

\begin{itemize}
\item Да се дадат определенията за контекстносвободна граматика, извод, език $L(G)$ и контекстносвободен език.
\item Да се обясни връзката между извод, най-ляв извод и дърво на синтактичен анализ.
\item Да се дадат конструкциите за затвореност относно обединение, конкатенация и звезда на Клини за \emph{произволни} входни граматики: преименуване на нетерминалите до непресичащи се множества, нов начален нетерминал и правилата $S\to S_1\mid S_2$, $S\to S_1S_2$, $S\to S_1S\mid\varepsilon$. Конспектът иска конструкциите, но не и доказателства към тях.
\item Да се формулира точно лемата за разрастването: разлагането $\alpha=xyuvw$, условията $|yv|\geq 1$, $|yuv|\leq p$ и $xy^iuv^iw\in L$ за всяко $i\in\mathbb{N}$.
\item Да се възпроизведе доказателствената идея чрез дърво в нормална форма на Чомски: двоично дърво, дълъг път, принцип на Дирихле и повторение на нетерминал.
\item Да се обясни защо повтореният нетерминал позволява „напомпване“ на частите $y$ и $v$.
\item Да се докаже с лемата, че $\{a^n b^n c^n\mid n\in\mathbb{N}\}$ не е КСЕ, като се разгледат всички допустими разлагания.
\item Да се дадат КСЕ $L_1=\{a^n b^n c^k\}$ и $L_2=\{a^n b^k c^k\}$, чието сечение е $\{a^n b^n c^n\}$.
\item Да се изведе незатвореността относно допълнение от незатвореността относно сечение чрез затвореността относно обединение и закона на де Морган.
\end{itemize}
